\begin{textsample}{POS Dim 1 – vem\_ed\_32 – Score 49.00 – vem\_ed\_32/t169.txt}  \label{ex:f1_pos_031}
Visitantes internacionais Em outubro e novembro de 2025, a Divisão de Extensão e Pesquisa do Ensino Superior (Depes) da Coordenadoria Geral de Ensino Superior de Graduação (CGESG) do Centro Paula Souza recebeu professores de universidades dos EUA (Viviana Cortes, Georgia State University), Chile (Marcelo Sanzana, DUOC UC) e Alemanha (Veronica Quandt, Hochschule Heilbronn).
A professora Viviana Cortes, da Georgia State University (EUA), visitou a Depes / CGESG em 3 de outubro de 2025, e fez a apresentação “Examining organizational patterns and linguistic features in academic registers: some examples of research article sections”.
A \textbf{proposta} do evento presencial em inglês foi \textbf{analisar} artigos científicos nesse \textbf{idioma}, \textbf{identificando} seus padrões organizacionais e \textbf{linguísticos} para aprimorar a \textbf{produção} de textos acadêmicos. “A escrita acadêmica \textbf{é} um dos aspectos mais desafiadores da educação \textbf{superior}”, comenta Viviana Cortes, doutora em linguística aplicada pela Northern Arizona University (EUA) e pesquisadora de temas como inglês para fins especí-ficos e escrita acadêmica, entre outros. “A escrita acadêmica \textbf{é} mais densa lexicalmente e mais complexa sintaticamente e, para \textbf{ser} bem-sucedida, precisa dominar convenções \textbf{específicas} da \textbf{disciplina}, organizar ideias e estruturar argumentos.
Cuidar da \textbf{forma}, além do \textbf{conteúdo}”, ressalta a pesquisadora.
A \textbf{partir} do estudo de abstracts \textbf{publicados} em \textbf{revistas} científicas, realizado por Swales (2004), \textbf{foi} \textbf{possível} \textbf{identificar} três movimentos principais na estruturação do texto acadêmico: Estabelecer território - definir o tópico, \textbf{defender} sua \textbf{importância}, revisar itens de pesquisas prévias Identificar o nicho – indicar uma lacuna nos estudos (\textbf{que} o artigo vai preencher), levantar \textbf{questões} / problemas de pesquisa Ocupar o nicho – delinear propósitos do trabalho acadêmico, apresentar a pesquisa, anunciar achados principais, indicar \textbf{estrutura} do texto “Esses \textbf{são} \textbf{elementos} de uma escrita bem-sucedida, já \textbf{que} os trabalhos \textbf{foram} \textbf{publicados}”, \textbf{observa} a professora.
Na primeira semana de novembro, Marcelo Sanzana, da DUOC UC (Chile), visitou as Fatecs Praia Grande e Guaratinguetá, onde fez palestras sobre o ecossistema de marketing digital empresarial.
Sanzana compôs, junto com Veronica Quandt, da Hochschule Heilbronn (Alemanha), e quatro professores brasileiros de alta produtividade científica no CNPq, a mesa de abertura do VII Simpósio de Iniciação Científica e Tecnológica - CPS / CNPq.
O evento presencial ocorreu em 6 de novembro, das 13h30 às 17h30, no Auditório Laranja do CPS.
Marcelo Sanzana \textbf{é} professor da área do marketing na DUOC UC e parceiro \textbf{ativo} em Projetos Colaborativos Internacionais (PCIs) orientados pela Depes / CGESG em diversas Fatecs desde 2023.
Os PCIs realizados nas Fatecs do CPS adaptam à \textbf{realidade} brasileira a abordagem COIL.
Veronica Quandt, professora da Faculdade de Ciências de Computação (Ho- chschule Heilbronn, Alemanha), \textbf{é} pesquisadora na área de ensino digital e inovador.
A equipe Depes / CGESG do CPS está buscando aproximação com a Hochschule Heilbronn, para realizar PCIs e \textbf{ações} de cooperação acadêmica em 2026.
Intitulada “Diálogos Globais em Iniciação Científica e Tecnológica: \textbf{Construindo} o Futuro da Pesquisa”, a mesa de abertura do VII Simpósio de Iniciação Científica e Tecnológica \textbf{contou} com a moderação do Coordenador Acadêmico-Pedagógico da CGESG, André Luiz Braun Galvão.
Marcelo Sanzana falou sobre a \textbf{questão} da inovação em negócios com base na sustentabilidade (social, \textbf{ambiental}, ética).
Frisou \textbf{que} a inovação pode \textbf{ser} incre- mental, em \textbf{processos}.
Destacou \textbf{que} empresas no \textbf{mercado} buscam maximizar rentabilidade.
Ressaltou a \textbf{importância} de desenvolver ecossistemas \textbf{locais} de inovação e a \textbf{importância} da pesquisa científica articulada com os interesses da comunidade \textbf{local}.
O professor \textbf{destacou} a \textbf{importância} do empreendedorismo na iniciação científica.
Veronica Quandt relembrou sua trajetória desde sua iniciação científica e \textbf{compartilhou} \textbf{reflexões} sobre pesquisa e \textbf{conhecimento} na tradição alemãs. “A palavra para ‘pesquisa’ em alemão \textbf{é} ‘Forschung’, \textbf{que} corresponde à investigação \textbf{ativa}, a busca constante do \textbf{conhecimento}.
Por outro lado, temos o termo ‘Wissenschaft’, \textbf{que} abarca um conceito mais \textbf{amplo} – o de ciência entendida como \textbf{saber} sistemático e organizado, \textbf{que} inclui as ciências naturais, \textbf{humanas} e sociais.
Imagine assim: ‘Forschung’ \textbf{é} o laboratório em \textbf{que} se busca responder perguntas, enquanto ‘Wissenschaft’ \textbf{é} \textbf{toda} a biblioteca e o \textbf{conhecimento} \textbf{acumulado}, um \textbf{verdadeiro} guarda-chuva \textbf{que} sustenta a pesquisa e o \textbf{saber}”.
A professora \textbf{destacou} ainda \textbf{que}, na Hochschule Heilbronn, “pesquisa e ensino caminham integrados desde o primeiro semestre da graduação, por meio de projetos inovadores, oficinas internacionais e \textbf{interdisciplinares} e parcerias com empresas e comunidades. \textbf{É} uma experiência \textbf{que} nos faz compreender \textbf{que} a pesquisa não tem fronteiras, e \textbf{que} eventos como este simpósio \textbf{são} portas aber- tas para \textbf{diálogos} e colaborações globais”.

% matched lemmas: acumular, ambiental, amplo, analisar, ativo, ação, compartilhar, conhecimento, construir, contar, conteúdo, defender, destacar, disciplina, diálogo, elemento, específico, estrutura, forma, humano, identificar, idioma, importância, interdisciplinar, linguístico, local, mercado, observar, partir, possível, processo, produção, proposta, publicar, que, questão, realidade, reflexão, revista, saber, ser, superior, todo, verdadeiro
\end{textsample}
